\chapter{Álgebra}

Aquí exploraremos los conceptos básicos
relacionados al álgebra. El álgebra
aparece en todas partes de la matemática,
y es importante conocer lo básico
ya que en particular se puede aplicar a
infinidad de problemas.

\section{Inducción}
Esto es un método fundamental en todas
las áreas de la matemática.

Para ilustrar qué dice lo podemos imaginar como
una serie de dominós en fila, tal que
si cae uno de ellos entonces cae el que
está inmediatamente a la derecha.
Es claro que si cae el primer dominó, entonces
caen todos los demás ya que si cae el primero,
entonces cae el segundo; si cae el segundo
entonces cae el tercero, y así sucesivamente.

% TODO: Añadir figura de dominós cayendo

Esta misma idea se puede formalizar, pero
en lugar de dominós, con proposiciones.
Esta es una definición informal:

\begin{theorem}[Inducción]
	\label{theorem:induccion}
	Si $P(n)$ es una proposición cuya veracidad
	depende del valor de $n$, se sabe que si
	$P(n)$ es cierto entonces $P(n + 1)$ es
	cierto, y $P(1)$ es cierto, entonces $P(n)$
	es cierto para todo $n$.
\end{theorem}

Para digerir el método lo mejor es usar
ejemplos.

\begin{example}[Suma Gaussiana]
	Demuestre que la suma de los primeros
	$n$ números naturales es igual a
	\[\frac{n(n + 1)}2.\]
\end{example}

Este resultado se le atribuye
al matemático Carl Gauss, quien, se cuenta,
lo dedujo cuando era un niño.

La manera de imaginar por qué
este resultado es cierto es sumar
los números de extremo a extremo como se
ve en la figura:

\begin{figure}[H]
	\begin{center}
		\input{figures/gaussian_sum.tex}
	\end{center}
	\caption{Suma Gaussiana.}
\end{figure}

Por el momento simplificamos el problema
y consideramos que $n$ es par. Claramente
la suma del primer con el último
número es $n + 1$,
del segundo con el penúltimo $2 + n - 1 = n + 1$,
y así sucesivamente todos los pares tienen la
misma suma. Pero nos damos cuenta de que
hay exactamente $\frac n2$ pares, de ahí la suma
necesariamente debe ser igual a
\[\frac n2 \cdot (n + 1).\]

Para demostrar con inducción que la fórmula
es cierta para todo $n$
podemos verificar que la misma funciona
para $n = 1$ ya que si sustituimos obtenemos
\[\frac{1 \cdot 2}{2} = 1\]
que es la suma hasta el primer número natural.

Ahora notemos que si conseguimos probar que
si la fórmula funciona para $n$, entonces
funciona para $n + 1$, entonces funcionaría
para todo $n$ ya que la misma funciona para
$n = 1$\footnote{Ten en cuenta el \autoref{theorem:induccion}}

Para ello si ocurre que
\[1 + 2 + \cdots + n = \frac{n(n + 1)}2,\]
debemos probar que
\[1 + 2 + \cdots + n + (n + 1) = \frac{(n + 1)(n + 2)}2,\]
pero sustituyendo
\[1 + 2 + \cdots + n = \frac{n(n + 1)}2\]
notamos que
\[\frac{n(n + 1)}2 + (n + 1) = (n + 1)\left(\frac n2 + 1\right) = \frac{(n + 1)(n + 2)}2,\]
que es lo que queríamos probar.

\subsection*{Problemas para esta sección}
\begin{problem}
Muestre que
\[1^2 + 3^2 + \cdots + n^2 = \frac{n(n + 1)(2n + 1)}6.\]
\end{problem}

\begin{problem}
Pruebe que la suma de los primeros
$n$ números impares es igual a $n^2$.
\end{problem}

\section{Funciones}
Algo recurrente son las funciones

\section{Desigualdades}

\section{Identidades}
Manejar expresiones algebraicas puede ser
bastante complicado si no se es capaz
de simplificar la misma.

E aquí algunos métodos y sustituciones
útiles.

\subsection{Trinomio cuadrado perfecto}
El trinomio cuadrado perfecto es quizás
el caso de factoreo más común de todos,
y tan solo consiste en saber que
\[(a + b)^2 = a^2 + 2ab + b^2.\]

En los problemas, a veces es conveniente
expandir la expresión $(a + b)^2$
y en otros casos resulta útil simplificar
$a^2 + 2ab + b^2$ como $(a + b)^2$.

\subsection{Super identidad}
Un amigo me contó que esta simple
transformación aparece en muchos lugares
inesperados y especialmente en teoría de números:
\[ab + cd = a(b - c) + c(d + a)\]

\subsection{Identidad de Sophie Germain}
La identidad de Sophie Germain establece lo siguiente:
\[a^4 + 4b^4 = (a^2 + 2b^2 + 2ab)(a^2 + 2b^2 - 2ab)\]

\subsection{Racionalización}
Normalmente es difícil tratar con raíces en
los denominadores de fracciones,
y el simple hecho de que
\[(a + b)(a - b) = a^2 - b^2\]
nos ayuda bastante ya que si
$a$ o $b$ son raíces cuadradas, es posible
multiplicar la expresión por $a - b$
para deshacerse de las raíces,
así, un uso común es cuando tenemos algo del
siguiente estilo:

\[\frac1{\sqrt a + \sqrt b}.\]

Podemos multiplicar la fracción
por $1$ y eso no cambia su valor,
pero
\[1 = \frac{\sqrt a - \sqrt b}{\sqrt a - \sqrt b}.\]

y si lo aplicamos a la expresión anterior vemos que
\[\frac1{\sqrt a + \sqrt b}\cdot\frac{\sqrt a - \sqrt b}{\sqrt a - \sqrt b} = \frac{\sqrt a - \sqrt b}{\sqrt a ^2 - \sqrt b ^2} = \frac{\sqrt a - \sqrt b}{a - b}.\]

Y ahora ya no tenemos raíces en el denominador.
Esta transformación es especialmente
últil si $a$ y $b$ son números
enteros o racionales, o
si $a - b$ lo es.

\section{Ecuaciones}

\section{Problemas}

\begin{problem}
Resuelva el sistema de ecuaciones:
\begin{equation*}
	\begin{cases}
		a^2 + b^2 + ab = 7 \\
		a - b          = 2.
	\end{cases}
\end{equation*}
\end{problem}

\begin{problem}
Si $a$ es un número real distinto de $0$, encuentre el valor de la expresión:
\[\frac{a + a^2 + a^3 + \cdots + a^n}{\frac{1}{a} + \frac{1}{a^2} + \frac{1}{a^3} + \cdots + \frac{1}{a^n}}\]
\end{problem}

\begin{problem}[Omapa 2023 - Ronda Departamental]
Encuentre el valor de la expresión
\[\frac2{\sqrt1 + \sqrt2} + \frac1{\sqrt2 + \sqrt3} + \cdots + \frac1{\sqrt{2022} + \sqrt{2023}}.\]
\, %%%% <----- WARNING: Si quitás el espacio la compilación crashea
\begin{hints}

	\begin{hint}
		Racionalización.
	\end{hint}
\end{hints}
\end{problem}

\begin{problem}[Libro ruso]
Demuestre que si los números positivos $a$, $b$ y $c$
forman una progresión aritmética,
entonces los números
\[\frac1{\sqrt b + \sqrt c}, \frac1{\sqrt c + \sqrt a}, \frac1{\sqrt a + \sqrt b}\]
también forma una progresión aritmética.
\end{problem}

\begin{problem}
Se disponen dos tazas, una con leche y otra con té,
y ambas tienen la misma cantidad de leche y de té.
A continuación se pone una cucharadita de la taza de
leche a la taza de té. Después se pone una cucharadita
de la taza de té en la taza de leche.
¿Qué hay más, leche en la taza de té o té en la taza de leche?
\end{problem}

\begin{problem}
Encuentre el valor de la expresión:
\[\frac1{1\cdot2} + \frac1{2\cdot3} + \frac1{3\cdot4} + \cdots + \frac1{n(n + 1)}.\]

% \begin{hints}
% 	\begin{hint}
% 		$\frac{1}{n} - \frac{1}{n + 1} = \frac1{n(n + 1 )}$.
% 	\end{hint}
% \end{hints}
\end{problem}

\begin{problem}
Se eligen $n$ números de la siguiente
figura, uno por cada fila y columna:

\begin{figure}[H]
	\begin{center}
		\input{figures/gridconstant.tex}
	\end{center}
\end{figure}

Encuentre todos los posibles valores
para la suma de todos los números seleccionados.
\end{problem}
